\chapter{ТЕОРЕТИЧЕСКИЕ ОСНОВЫ РАЗРАБОТКИ ВЕБ-ПРИЛОЖЕНИЯ ВИЗУАЛИЗАЦИИ ИСТОРИЧЕСКИХ ВОЕННЫХ КОНФЛИКТОВ}

\section{Особенности визуализации данных об исторических военных событиях}

Предметная область в контексте данной работы представляет собой совокупность данных, методов и процессов, связанных с фиксацией, анализом и представлением информации о ходе военных конфликтов.

С технической точки зрения ключевой задачей является систематизация и цифровизация разнородных сведений, включающих текстовые описания, карты, табличные данные и хроники, а также их преобразование в формализованные сущности, пригодные для хранения в базе данных, последующей обработки и визуализации.

Ключевым объектом исследования выступают боевые действия оперативно-стратегического уровня. Это предполагает фокус на масштабе армейских операций и сражений, основных группировках войск, динамике линии фронта и перемещениях соединений, таких как дивизии, корпуса и армии. Основной задачей визуализации на этом уровне является отражение не статичной картины сражения, а его пространственно-временной динамики. Для этого данные должны быть представлены в машиночитаемом виде с явной привязкой к координатам карты и времени, а также в согласованных форматах обмена. В рамках разрабатываемой системы основными форматами такого обмена целесообразно рассматривать JSON для справочной и атрибутивной информации и GeoJSON для пространственных векторных объектов~\cite{rfc7946}. Форматы KML и Shapefile могут использоваться как внешние форматы импорта или экспорта геоданных, а сервисы WMTS --- для подключения готовых картографических слоёв и подложек~\cite{ogc_wmts}.

Основные сущности предметной области включают:
\begin{itemize}
\item стороны конфликта, представляющие собой государства, коалиции и военно-политические блоки; каждая сторона обладает своим составом войск, командной структурой и стратегическими целями;
\item войсковые соединения и объединения, такие как армии, корпуса и дивизии, являющиеся основными оперативными единицами на карте; их существенными атрибутами выступают тип и принадлежность;
\item линию фронта и зоны контроля, представляемые динамически изменяющимися полигонами, отражающими территорию, контролируемую сторонами конфликта в конкретный момент времени;
\item боевые операции и сражения, представляющие собой ограниченные во времени и пространстве периоды активных боевых действий, имеющие собственное название, даты проведения, цели и итоги;
\item временную шкалу, с помощью которой возможно прослеживать развитие событий в ходе боевых операций.
\end{itemize}

Для данных, изменяющихся во времени, существенное значение имеет выбор способа временной навигации, поскольку пользователь должен иметь возможность сопоставлять состояние объектов в разные моменты и отслеживать развитие процесса~\cite{aigner_time_oriented_data}.

Традиционные методы представления подобной информации, прежде всего статичные карты в учебниках и атласах, имеют принципиальные ограничения. Они вынуждены либо показывать ситуацию на одну фиксированную дату, либо использовать множество стрелок и условных обозначений на одном листе, что затрудняет восприятие общей динамики. Следовательно, предметная область нуждается в переходе от статичных моделей к интерактивным, управляемым данными средствам визуализации, где время выступает координатой для навигации и анализа.

При проектировании картографического интерфейса важно учитывать не только точность отображаемых данных, но и особенности восприятия карты пользователем, поскольку выбор условных обозначений и способов представления пространственных взаимосвязей влияет на понимание отображаемых событий~\cite{maceachren_how_maps_work}.

Такой подход позволяет преобразовать линейное текстовое или статичное графическое описание конфликта в единую интерактивную модель, в рамках которой исследователь может:
\begin{itemize}
\item наблюдать изменение линии фронта в ходе проведения военных операций;
\item синхронизировать события на различных театрах военных действий;
\item выявлять взаимосвязь между перемещениями соединений и ключевыми сражениями.
\end{itemize}

Таким образом, анализ предметной области подтверждает актуальность разработки системы «Палантир» и выявляет ключевые требования к ней: необходимость систематизации и хранения данных в структурированном виде, работы с динамическими пространственно-временными данными, обеспечения интерактивности и учёта специфики оперативно-стратегического уровня ведения боевых действий. При этом уже существуют системы и проекты, которые в той или иной степени реализуют отдельные элементы указанных задач, прежде всего картографическую привязку и оперативное обновление данных, что делает целесообразным их дальнейший обзор и сравнительный анализ.

\section{Анализ существующих программных решений и аналогичных проектов}

Для формирования обоснованной концепции веб-приложения «Палантир» был проведён сравнительный анализ существующих цифровых платформ, специализирующихся на визуализации военных конфликтов. Анализ выявил несколько устойчивых категорий решений, каждая из которых обладает своими преимуществами и ограничениями.

\textbf{Специализированные платформы для мониторинга текущих конфликтов.} К данной категории относятся такие ресурсы, как Lostarmour, DivGen и Military Summary. Эти проекты демонстрируют высокую оперативность обновления данных и детализацию на тактическом уровне, часто основанную на анализе открытых источников. Сводные характеристики указанных решений приведены в таблице~\ref{tab:specialized-platforms}.

Преимуществами таких платформ являются:
\begin{itemize}
\item высокая детализация и актуальность информации;
\item активное использование географических привязок и картографической основы;
\item ориентация на отображение динамически меняющейся линии соприкосновения.
\end{itemize}

К их недостаткам следует отнести:
\begin{itemize}
\item ориентацию преимущественно на один текущий конфликт без охвата других исторических периодов;
\item отсутствие исторической справки, предпосылок и итогов операций;
\item риск неполноты или недостаточной верифицированности данных в условиях высокой скорости обновления.
\end{itemize}

\begin{table}[H]
\caption{\label{tab:specialized-platforms}Характеристики специализированных платформ мониторинга текущих конфликтов}
\centering
{\footnotesize
\begin{tabular}{|p{2.2cm}|p{4.2cm}|p{4.2cm}|p{4.2cm}|}
\hline
\textbf{Критерий} & \textbf{Lostarmour} & \textbf{DivGen} & \textbf{Military Summary} \\ \hline
Основная функция & Предоставление интерактивной карты и аналитических материалов по текущему конфликту & Предоставление интерактивной карты оперативной обстановки и кратких статей о воинских подразделениях & Предоставление интерактивной карты и кратких сводок об изменениях оперативной обстановки \\ \hline
Картографичес-кая основа & Реализована на базе Конструктора карт Яндекса & Используются картографические сервисы Яндекс и Google & Используются картографические сервисы Яндекс, Google и OpenStreetMap \\ \hline
Временная навигация & Реализована временная шкала & Реализована временная шкала & Реализована временная шкала \\ \hline
Детализация объектов и событий & Высокий уровень детализации: отдельные эпизоды, подсчёт потерь единиц техники, статистические данные и изменение линии фронта & Высокий уровень детализации: отдельные эпизоды, описание и размещение подразделений, изменения линии фронта & Средний уровень детализации: события новостного характера и изменения линии фронта \\ \hline
Поиск и фильтрация & Ограниченные возможности фильтрации; навигация по тематическим разделам и материалам & Ограниченные возможности фильтрации; базовая навигация по карте и участкам & Ограниченные возможности фильтрации; навигация по карте и текстовым сводкам \\ \hline
Информацион-ный контекст & Развёрнутые пояснительные и аналитические материалы в рамках одного конфликта & Краткие пояснительные и аналитические материалы & Краткие сводные материалы с минимальным аналитическим сопровождением \\ \hline
\end{tabular}}
\end{table}

\textbf{Универсальные картографические сервисы.} К данной категории относятся Google My Maps, Конструктор карт Яндекса и другие сервисы на базе OpenStreetMap. Они предоставляют пользователям базовые инструменты для создания собственных карт с маркерами, линиями и фигурами. Сравнение функциональных возможностей наиболее распространённых сервисов представлено в таблице~\ref{tab:map-services}.

\begin{table}[H]
\caption{\label{tab:map-services}Сравнение Google My Maps и Конструктора карт Яндекса}
\centering
{\footnotesize
\begin{tabular}{|p{3.0cm}|p{6.3cm}|p{5.7cm}|}
\hline
\textbf{Критерий} & \textbf{Google My Maps} & \textbf{Конструктор карт Яндекса} \\ \hline
Идентификация и авторизация & Аккаунт Google & Аккаунт Яндекс \\ \hline
Данные и импорт & Импорт данных в форматах CSV, XLSX, GeoJSON, KML/KMZ, GPX; привязка по столбцам адресов и координат & Импорт данных в форматах CSV, Excel, GeoJSON и KML \\ \hline
Слои и группировка & Возможно создание нескольких слоёв и группировка объектов по слоям & Отсутствуют полноценные слои и группировка объектов \\ \hline
Совместная работа & Совместное редактирование и публикация & Публикация и встраивание; совместное редактирование ограничено \\ \hline
Встраивание на сайт & Встраивание через iframe-код & Встраивание через iframe-код \\ \hline
Стиль подложки & Базовые типы подложки Google без глубокой кастомизации & Базовые типы подложки Яндекс.Карт без глубокой кастомизации \\ \hline
\end{tabular}}
\end{table}

Преимуществами универсальных картографических сервисов являются:
\begin{itemize}
\item доступность и простота создания несложных карт;
\item проработанная и актуальная картографическая подложка.
\end{itemize}

Недостатки таких сервисов заключаются в следующем:
\begin{itemize}
\item отсутствует возможность привязки объектов на карте к временной шкале и анимации их изменений, что является ключевым ограничением;
\item отсутствуют средства для отображения воинских соединений, типов операций, линий фронта и зон контроля;
\item информация представляется в виде произвольных меток, а не единой связанной базы данных, что затрудняет сложный поиск и анализ.
\end{itemize}

\textbf{Академические и энциклопедические исторические проекты.} К данной категории можно отнести интерактивные карты на образовательных порталах и проекты, подобные «The Fallen of World War II». Такие решения сопровождаются поясняющими текстами, инфографикой и визуальными сценариями.

Их преимуществами являются:
\begin{itemize}
\item высокая достоверность и научная обоснованность данных;
\item качественная визуализация и продуманный нарратив.
\end{itemize}

В качестве недостатков можно выделить:
\begin{itemize}
\item отсутствие полноценного взаимодействия пользователя с данными;
\item ограниченную тематику, поскольку такие проекты часто посвящены одному конкретному конфликту или отдельному его аспекту.
\end{itemize}

\textbf{Открытые wiki-платформы.} Речь идёт о проектах, построенных на принципах вики, где контент создаётся сообществом энтузиастов.

Их достоинствами являются:
\begin{itemize}
\item большой охват и объём накопленной информации;
\item возможность коллективной работы над данными.
\end{itemize}

Основные недостатки таких платформ состоят в том, что:
\begin{itemize}
\item информация часто хранится в виде неструктурированного текста, а не формализованных записей, что делает её малопригодной для автоматизированной визуализации;
\item картографическая составляющая выражена слабо, а карты, если и присутствуют, чаще являются статичными изображениями, а не интерактивными системами;
\item из-за открытой модели редактирования данные могут быть противоречивыми и требуют постоянной строгой модерации.
\end{itemize}

Сравнительный анализ позволил выявить пробел в классе подобных решений. Существующие проекты либо обеспечивают детализацию в реальном времени без исторического контекста и аналитического инструментария, либо представляют универсальные инструменты без военной специфики, либо являются закрытыми образовательными продуктами с ограниченными возможностями настройки, либо выступают в роли неструктурированных баз знаний без развитой визуализации.

\section{Концептуальная модель веб-приложения визуализации данных о военных действиях «Палантир»}

На основании проведённого сравнительного анализа существующих решений была разработана концептуальная модель системы «Палантир». Цель проекта заключается в заполнении выявленных пробелов за счёт создания комплекса средств визуализации и анализа. Система предназначена для представления данных о военных конфликтах с привязкой к месту и времени.

Концептуальная архитектура построена вокруг трёх ключевых компонентов:
\begin{itemize}
\item единой базы данных структурированной информации, содержащей нормализованные сведения о конфликтах, сторонах, соединениях и изменениях линии фронта с привязкой к координатам и времени;
\item пользовательского интерфейса, сочетающего интерактивную карту с элементом управления временной шкалой и позволяющего осуществлять навигацию во времени;
\item модуля визуализации и анализа, отвечающего за отображение данных на карте в виде слоёв и предоставление инструментов фильтрации и поиска.
\end{itemize}

Такое разделение соответствует рекомендуемой многоуровневой архитектуре веб-ориентированных геоинформационных систем, обеспечивающей производительность, масштабируемость и надёжность за счёт чёткого разграничения уровней хранения данных, логики и представления~\cite{sommerville2015software}.

Для эффективной работы с динамическими пространственно-временными данными, такими как перемещение соединений и изменение линии фронта, архитектура приложения должна опираться на современные подходы к обработке информации. Это позволяет снизить задержки при визуализации и обеспечить стабильную производительность при работе с большими объёмами данных.

Функциональные требования к приложению включают:
\begin{itemize}
\item линейную временную шкалу с возможностью выбора конкретной даты, обеспечивающую сопоставление состояний объектов в разные моменты времени~\cite{aigner_time_oriented_data};
\item избирательное отображение слоёв, включая войсковые соединения, линию фронта и зоны контроля;
\item возможность использования различных карт-подложек, таких как спутниковые снимки, исторические и схематичные карты;
\item отображение дислокации воинских соединений с использованием условных обозначений;
\item визуализацию динамически изменяющихся полигонов, отображающих зоны контроля сторон конфликта;
\item информационные карточки, содержащие сведения о событиях, происходивших в конкретном месте и в конкретное время;
\item фильтрацию отображаемых данных по стороне конфликта, типу соединения и периоду времени.
\end{itemize}

К нефункциональным требованиям относятся:
\begin{itemize}
\item расширяемость архитектуры, позволяющая добавлять новые конфликты и исторические периоды за счёт расширения базы данных без изменения базовой логики приложения;
\item возможность обработки значительных объёмов данных и поддержки работы с пространственными векторными объектами;
\item структурированность данных, обеспечивающая однозначную интерпретацию и пригодность сведений для машинной обработки;
\item доступность приложения без установки дополнительного специализированного программного обеспечения;
\item удобство использования, предполагающее наличие навигации по карте и временной шкале;
\item совместимость с открытыми стандартами, такими как GeoJSON и Web Map Service, позволяющую подключать внешние картографические источники и использовать географически привязанные изображения карт~\cite{rfc7946,ogc_wms}.
\end{itemize}

Концепция системы «Палантир» заключается в создании не только картографического сервиса, но и интерактивной среды для освещения хода боевых действий. В отличие от статичных карт, система позволяет рассматривать боевые действия как пространственно-временной процесс, в котором военные операции противоборствующих сторон взаимосвязаны. В отличие от платформ мониторинга текущих событий, она предоставляет исторический контекст и возможность сравнения различных конфликтов.

\section{Обоснование выбора технологического стека}

Выбор технологического стека для системы «Палантир» определяется характером обрабатываемых данных и требованиями к их представлению. Система ориентирована на хранение и выдачу взаимосвязанных сущностей предметной области, включая конфликты, стороны, операции и соединения, а также на работу с пространственными объектами и временными характеристиками. Дополнительно требуется обеспечить сетевое взаимодействие клиентской и серверной частей по HTTP и передачу данных в распространённых форматах обмена. В рамках обоснования рассматриваются язык программирования для серверной части, веб-платформа для построения API, система управления базами данных с поддержкой пространственных типов, библиотека Leaflet, предоставляющая во фронтенде средства создания интерактивной карты, работы со слоями, маркерами, полилиниями и полигонами~\cite{leaflet_api_reference}, а также набор технологий платформы .NET, применяемых при реализации.

\subsection{Выбор языка программирования и серверной платформы ASP.NET Core}

Для реализации серверной части системы выбран язык программирования C\#, функционирующий в рамках платформы .NET. Такой выбор обусловлен тем, что серверная логика предполагает обработку структурированных данных, содержащих устойчивые связи между сущностями и наборами атрибутов, а также выполнение операций выборки по фильтрам и временным параметрам. Статическая типизация C\# обеспечивает контроль структуры данных на этапе компиляции и позволяет формализовать предметную область в виде классов, перечислений и связей между объектами. Подобная модель используется для описания сущностей и передачи данных между слоями приложения. Применение объектно-ориентированного подхода позволяет соотнести программную модель с концептуальной моделью системы и обеспечить однозначное представление основных объектов предметной области.

Важным для серверных приложений является наличие средств для построения запросов к данным и обработки коллекций. В рамках платформы .NET используются механизмы обобщений и язык интегрированных запросов LINQ, применяемые при формировании выборок, фильтрации и преобразовании результатов для последующей выдачи через API. Для серверной части также значимо использование асинхронной модели выполнения, позволяющей выполнять операции ввода-вывода без блокирования потока исполнения. Это особенно важно при обслуживании параллельных обращений к API и при выполнении запросов к базе данных~\cite{dotnet_docs,linq}.

Таким образом, применение C\# и платформы .NET соответствует требованиям к реализации серверной части, ориентированной на работу со структурированными и пространственно-временными данными.

В качестве технологии разработки серверной части выбрана платформа ASP.NET Core. Система предполагает взаимодействие клиентского приложения с сервером посредством HTTP-запросов. ASP.NET Core ориентирована на создание веб-приложений и REST-сервисов, в том числе с передачей данных в формате JSON. Платформа предоставляет механизмы маршрутизации запросов, обработки параметров, связывания входных данных с моделями и формирования ответов. Это используется для построения API, обеспечивающего доступ к данным системы и поддержку прикладных сценариев, например выборку данных на дату или интервал, получение списков сущностей и пространственных объектов для отображения на карте. ASP.NET Core поддерживает модульную организацию приложения за счёт конвейера промежуточного программного обеспечения, применяемого для обработки ошибок, настройки политик доступа, ведения журналирования и работы с конфигурацией. Встроенный механизм внедрения зависимостей используется для изоляции компонентов и управления их взаимодействием, включая подключение репозиториев, сервисов прикладного уровня и контекстов доступа к данным. Следовательно, использование ASP.NET Core обеспечивает реализацию серверного API, соответствующего выбранной модели сетевого взаимодействия и требованиям к структурированию серверной части~\cite{aspnet_docs}.

\subsection{Выбор системы управления базами данных PostgreSQL/PostGIS}

Разрабатываемая система предполагает хранение структурированных данных и пространственных объектов, отражающих линию фронта, зоны контроля и расположение соединений. В качестве системы управления базами данных выбрана PostgreSQL с расширением PostGIS. PostgreSQL обеспечивает поддержку реляционной модели данных, что позволяет реализовать нормализованную структуру хранения информации о конфликтах, сторонах, операциях, соединениях и временных характеристиках. Реляционная модель применяется для обеспечения целостности данных и поддержания связей между сущностями за счёт первичных и внешних ключей, ограничений целостности и транзакционных операций. Расширение PostGIS добавляет поддержку пространственных типов данных, включая точки, линии и полигоны. Это позволяет хранить геометрические объекты непосредственно в базе данных и выполнять пространственные операции, применимые к задачам визуализации, например проверку пересечений, принадлежности области и вычисление расстояний. Использование пространственных индексов обеспечивает выполнение пространственных запросов на стороне СУБД, что упрощает серверную обработку и снижает необходимость ручного анализа геометрий в прикладном коде. Для обмена пространственными данными с клиентской частью может использоваться формат GeoJSON, поддерживающий передачу геометрии и атрибутов в единой структуре. Такой подход применяется при подготовке данных для картографической визуализации в веб-интерфейсе~\cite{postgresql_docs,postgis_manual}.

Таким образом, выбор PostgreSQL с расширением PostGIS обусловлен необходимостью совместного хранения атрибутивной информации и пространственных данных в рамках единой СУБД, а также поддержкой операций, требуемых для пространственно-временной визуализации.

\subsection{Технологии платформы .NET, применяемые при разработке}

Для взаимодействия серверной части с базой данных используется Entity Framework Core — объектно-реляционный механизм отображения данных. Применение ORM-подхода позволяет сопоставлять таблицы базы данных с классами C\#, что обеспечивает согласованность структуры программной модели и структуры хранения. Формирование запросов осуществляется с использованием LINQ, интегрированного в язык программирования, при этом результат запросов может преобразовываться в структуры передачи данных, используемые API. Entity Framework Core поддерживает механизм миграций, позволяющий фиксировать изменения схемы базы данных в виде последовательных версий и применять их при развёртывании и обновлении. Такой механизм используется для синхронизации структуры базы данных между средами разработки и эксплуатации, а также для контролируемого изменения структуры при развитии модели данных~\cite{efcore_overview}.

Для построения серверного приложения применяются средства платформы .NET, связанные с обработкой сетевых запросов и сериализацией данных. В частности, используются механизмы конфигурации приложения, журналирования и внедрения зависимостей, обеспечивающие связывание компонентов приложения, управление их жизненным циклом и последовательностью обработки данных. Эти средства применяются при подключении контекста доступа к данным, регистрации сервисов прикладного уровня и настройке поведения API.

В качестве среды разработки используется Visual Studio, предоставляющая инструменты отладки, управления зависимостями через NuGet и интеграцию с системой контроля версий. Использование интегрированной среды разработки обеспечивает централизованное управление проектом и его компонентами, а также контроль конфигураций сборки и запуска.

Таким образом, выбранные технологии платформы .NET формируют основу серверной реализации и обеспечивают типовые механизмы доступа к данным, построения API и сопровождения разработки.
